
\documentclass{article}
\usepackage{colortbl}
\usepackage{makecell}
\usepackage{multirow}
\usepackage{supertabular}

\begin{document}

\newcounter{utterance}

\centering \large Interaction Transcript for game `clean\_up', experiment `2\_hard\_7obj\_ar', episode 2 with gemma{-}4{-}26b{-}a4b{-}it.
\vspace{24pt}

{ \footnotesize  \setcounter{utterance}{1}
\setlength{\tabcolsep}{0pt}
\begin{supertabular}{c@{$\;$}|p{.15\linewidth}@{}p{.15\linewidth}p{.15\linewidth}p{.15\linewidth}p{.15\linewidth}p{.15\linewidth}}
    \# & $\;$A & \multicolumn{4}{c}{Game Master} & $\;\:$B\\
    \hline

    \theutterance \stepcounter{utterance}  
    & & \multicolumn{4}{p{0.6\linewidth}}{
        \cellcolor[rgb]{0.9,0.9,0.9}{
            \makecell[{{p{\linewidth}}}]{
                \texttt{\tiny{[P1$\langle$GM]}}
                \texttt{أنا مدير لعبتك، وأنت تلعب لعبة تعاونية مع الشبكة التالية كلوحة لعب:} \\
\\ 
\texttt{```} \\
\texttt{ 1234567} \\
\texttt{╔═╤══╤══╗} \\
\texttt{║◌│◌T│◌◌║ 1} \\
\texttt{║◌└┬─┴──╢ 2} \\
\texttt{║◌◌│O◌◌◌║ 3} \\
\texttt{╟─┬┘P◌┌─╢ 4} \\
\texttt{║◌│AD◌│◌║ 5} \\
\texttt{╟─┼─┬─┤M║ 6} \\
\texttt{║◌│S│◌│◌║ 7} \\
\texttt{╚═╧═╧═╧═╝} \\
\\ 
\texttt{```} \\
\\ 
\texttt{* تعرض الحافة العلوية إحداثيات x التي تتزايد نحو اليمين، والحافة اليمنى إحداثيات y التي تتزايد نحو الأسفل.} \\
\texttt{* الكائنات التالية موضوعة عشوائياً على شبكتك: 'P', 'O', 'T', 'S', 'D', 'A', 'M'.} \\
\\ 
\\ 
\texttt{يرى اللاعب الآخر نسخة مختلفة من لوحة اللعب، حيث توضع الكائنات في مواقع عشوائية مختلفة. لا يمكنك رؤية لوحة اللاعب الآخر، ولا يمكنه رؤية لوحتك.} \\
\\ 
\\ 
\texttt{**الهدف:**} \\
\\ 
\texttt{يجب على كلا اللاعبين تحريك الكائنات على خلفية كل منهما بحيث تنتهي الكائنات المتطابقة عند الإحداثيات نفسها. عليك التواصل مع اللاعب الآخر للاتفاق على إعداد هدف مشترك.} \\
\\ 
\\ 
\texttt{**القواعد:**} \\
\\ 
\texttt{* في كل دور، يمكنك إرسال أمر واحد فقط من الأمرين التاليين:} \\
\texttt{ 1. `قل: <رسالة>`: لإرسال رسالة (كل شيء حتى فاصل السطر التالي) إلى اللاعب الآخر. سأقوم بتمريرها إلى شريكك.} \\
\texttt{2. `حرّك: <كائن>, (<X>, <Y>)`: لتحريك كائن إلى موضع جديد، حيث `<X>` هو العمود و `<Y>` هو الصف. سأخبرك بما إذا كانت حركتك ناجحة أم لا.} \\
\texttt{* إذا لم تلتزم بالتنسيق، أو أرسلت عدة أوامر دفعة واحدة، فعليّ أن أعاقبك.} \\
\texttt{* إذا تجاوز كلا اللاعبين أكثر من 16 عقوبة، فستخسران اللعبة معاً.} \\
\texttt{* من الضروري أن تتواصل مع اللاعب الآخر بشأن حالة هدفك! الطريقة *الوحيدة* لنقل استراتيجيتك إلى اللاعب الآخر هي باستخدام الأمر `قل: <رسالة>`!} \\
\\ 
\\ 
\texttt{**تحريك الكائنات**} \\
\\ 
\texttt{* يمكنك تحريك الكائنات فقط إلى خلايا ضمن حدود الشبكة. يجب أن تكون الخلية المستهدفة فارغة، أي يجب أن تحتوي فقط على الرمز '◌'.} \\
\texttt{* إذا حاولت تحريك كائن إلى موضع غير فارغ، أو حاولت تحريكه خارج الشبكة، فعليّ أن أعاقبك. ستحصل على محاولة أخرى.} \\
\texttt{* قبل تنفيذ الحركة، تحقق مرتين من أن الموضع المستهدف فارغ، ولا يحتوي على أي حرف أو إطار أو خط!} \\
\\ 
\\ 
\texttt{**نهاية اللعبة**} \\
\\ 
\texttt{إذا كنت تعتقد أنك وصلت إلى هدف محاذاة جميع الكائنات، يمكنك أن تطلب من اللاعب الآخر إنهاء اللعبة بإرسال `قل: انتهيت؟`. إذا طلب منك اللاعب الآخر إنهاء اللعبة، ورددت بـ `قل: انتهيت!`، ستنتهي اللعبة.} \\
\\ 
\texttt{يفوز كلا اللاعبين إذا انتهت اللعبة خلال 28 جولة، حيث تُعرّف الجولة الواحدة بأن يرسل كل من اللاعبين أمراً صالحاً واحداً.} \\
\\ 
\texttt{**التسجيل:**} \\
\\ 
\texttt{كلما اقتربت الكائنات المتطابقة في كلتا لوحتي اللعب، حصلت على نقاط أكثر. العقوبات تقلل نقاطك. هل يمكنك تحطيم الرقم القياسي؟} \\
\\ 
\\ 
\texttt{يرجى إرسال رسالة إلى اللاعب الآخر لبدء اللعبة!} \\
            }
        }
    }
    & \\ \\

    \theutterance \stepcounter{utterance}  
    & \multicolumn{4}{p{0.6\linewidth}}{
        \cellcolor[rgb]{0.9,0.9,0.9}{
            \makecell[{{p{\linewidth}}}]{
                \texttt{\tiny{[P1$\rangle$GM]}}
                \texttt{قل: مرحباً! لنبدأ بالتعاون. سأقوم أولاً بمشاركة مواقع الكائنات الموجودة في لوحتي لنحاول إيجاد هدف مشترك.} \\
\\ 
\texttt{مواقع الكائنات لدي هي:} \\
\texttt{T: (3, 1)} \\
\texttt{O: (3, 3)} \\
\texttt{P: (4, 4)} \\
\texttt{A: (2, 5)} \\
\texttt{D: (3, 5)} \\
\texttt{S: (2, 7)} \\
\texttt{M: (7, 6)} \\
\\ 
\texttt{ما هي مواقع الكائنات لديك؟} \\
            }
        }
    }
    & & \\ \\

    \theutterance \stepcounter{utterance}  
    & & & \multicolumn{2}{p{0.3\linewidth}}{
        \cellcolor[rgb]{0.9,0.9,0.9}{
            \makecell[{{p{\linewidth}}}]{
                \texttt{\tiny{[GM$|$GM]}}
                \texttt{قل: مرحباً! لنبدأ بالتعاون. سأقوم أولاً بمشاركة مواقع الكائنات الموجودة في لوحتي لنحاول إيجاد هدف مشترك.} \\
\\ 
\texttt{مواقع الكائنات لدي هي:} \\
\texttt{T: (3, 1)} \\
\texttt{O: (3, 3)} \\
\texttt{P: (4, 4)} \\
\texttt{A: (2, 5)} \\
\texttt{D: (3, 5)} \\
\texttt{S: (2, 7)} \\
\texttt{M: (7, 6)} \\
\\ 
\texttt{ما هي مواقع الكائنات لديك؟} \\
            }
        }
    }
    & & \\ \\

    \theutterance \stepcounter{utterance}  
    & & & \multicolumn{2}{p{0.3\linewidth}}{
        \cellcolor[rgb]{0.9,0.9,0.9}{
            \makecell[{{p{\linewidth}}}]{
                \texttt{\tiny{[GM$|$GM]}}
                \texttt{Invalid format: tail is not empty: '} \\
\\ 
\texttt{مواقع الكائنات لدي هي:} \\
\texttt{T: (3, 1)} \\
\texttt{O: (3, 3)} \\
\texttt{P: (4, 4)} \\
\texttt{A: (2, 5)} \\
\texttt{D: (3, 5)} \\
\texttt{S: (2, 7)} \\
\texttt{M: (7, 6)} \\
\\ 
\texttt{ما هي مواقع الكائنات لديك؟} \\
            }
        }
    }
    & & \\ \\

    \theutterance \stepcounter{utterance}  
    & & \multicolumn{4}{p{0.6\linewidth}}{
        \cellcolor[rgb]{0.9,0.9,0.9}{
            \makecell[{{p{\linewidth}}}]{
                \texttt{\tiny{[P1$\langle$GM]}}
                \texttt{عقوبة: يجب ألا تحتوي رسالتك على أي شيء بعد الأمر!} \\
\\ 
\texttt{تأكد من أن ردك يحتوي فقط على `قل: <رسالة>` أو `حرّك: <كائن>, (<X>, <Y>)`، ولا شيء غير ذلك!} \\
\\ 
\texttt{لقد جمعتما معاً 1 من أصل 16 عقوبة.} \\
\texttt{يرجى المحاولة مرة أخرى!} \\
            }
        }
    }
    & \\ \\

    \theutterance \stepcounter{utterance}  
    & \multicolumn{4}{p{0.6\linewidth}}{
        \cellcolor[rgb]{0.9,0.9,0.9}{
            \makecell[{{p{\linewidth}}}]{
                \texttt{\tiny{[P1$\rangle$GM]}}
                \texttt{قل: مرحباً! لنبدأ بالتعاون. سأشارك مواقع الكائنات في لوحتي لنحدد هدفاً مشتركاً.} \\
\\ 
\texttt{T: (3, 1)} \\
\texttt{O: (3, 3)} \\
\texttt{P: (4, 4)} \\
\texttt{A: (2, 5)} \\
\texttt{D: (3, 5)} \\
\texttt{S: (2, 7)} \\
\texttt{M: (7, 6)} \\
\\ 
\texttt{ما هي مواقع الكائنات لديك؟} \\
            }
        }
    }
    & & \\ \\

    \theutterance \stepcounter{utterance}  
    & & & \multicolumn{2}{p{0.3\linewidth}}{
        \cellcolor[rgb]{0.9,0.9,0.9}{
            \makecell[{{p{\linewidth}}}]{
                \texttt{\tiny{[GM$|$GM]}}
                \texttt{قل: مرحباً! لنبدأ بالتعاون. سأشارك مواقع الكائنات في لوحتي لنحدد هدفاً مشتركاً.} \\
\\ 
\texttt{T: (3, 1)} \\
\texttt{O: (3, 3)} \\
\texttt{P: (4, 4)} \\
\texttt{A: (2, 5)} \\
\texttt{D: (3, 5)} \\
\texttt{S: (2, 7)} \\
\texttt{M: (7, 6)} \\
\\ 
\texttt{ما هي مواقع الكائنات لديك؟} \\
            }
        }
    }
    & & \\ \\

    \theutterance \stepcounter{utterance}  
    & & & \multicolumn{2}{p{0.3\linewidth}}{
        \cellcolor[rgb]{0.9,0.9,0.9}{
            \makecell[{{p{\linewidth}}}]{
                \texttt{\tiny{[GM$|$GM]}}
                \texttt{Invalid format: tail is not empty: '} \\
\\ 
\texttt{T: (3, 1)} \\
\texttt{O: (3, 3)} \\
\texttt{P: (4, 4)} \\
\texttt{A: (2, 5)} \\
\texttt{D: (3, 5)} \\
\texttt{S: (2, 7)} \\
\texttt{M: (7, 6)} \\
\\ 
\texttt{ما هي مواقع الكائنات لديك؟} \\
            }
        }
    }
    & & \\ \\

    \theutterance \stepcounter{utterance}  
    & & \multicolumn{4}{p{0.6\linewidth}}{
        \cellcolor[rgb]{0.9,0.9,0.9}{
            \makecell[{{p{\linewidth}}}]{
                \texttt{\tiny{[P1$\langle$GM]}}
                \texttt{عقوبة: يجب ألا تحتوي رسالتك على أي شيء بعد الأمر!} \\
\\ 
\texttt{تأكد من أن ردك يحتوي فقط على `قل: <رسالة>` أو `حرّك: <كائن>, (<X>, <Y>)`، ولا شيء غير ذلك!} \\
\\ 
\texttt{لقد جمعتما معاً 2 من أصل 16 عقوبة.} \\
\texttt{يرجى المحاولة مرة أخرى!} \\
            }
        }
    }
    & \\ \\

    \theutterance \stepcounter{utterance}  
    & \multicolumn{4}{p{0.6\linewidth}}{
        \cellcolor[rgb]{0.9,0.9,0.9}{
            \makecell[{{p{\linewidth}}}]{
                \texttt{\tiny{[P1$\rangle$GM]}}
                \texttt{قل: مرحباً! لنبدأ بالتعاون. سأشارك مواقع الكائنات في لوحتي لنحدد هدفاً مشتركاً. T:(3,1), O:(3,3), P:(4,4), A:(2,5), D:(3,5), S:(2,7), M:(7,6). ما هي مواقع الكائنات لديك؟} \\
            }
        }
    }
    & & \\ \\

    \theutterance \stepcounter{utterance}  
    & & & \multicolumn{2}{p{0.3\linewidth}}{
        \cellcolor[rgb]{0.9,0.9,0.9}{
            \makecell[{{p{\linewidth}}}]{
                \texttt{\tiny{[GM$|$GM]}}
                \texttt{قل: مرحباً! لنبدأ بالتعاون. سأشارك مواقع الكائنات في لوحتي لنحدد هدفاً مشتركاً. T:(3,1), O:(3,3), P:(4,4), A:(2,5), D:(3,5), S:(2,7), M:(7,6). ما هي مواقع الكائنات لديك؟} \\
            }
        }
    }
    & & \\ \\

    \theutterance \stepcounter{utterance}  
    & & \multicolumn{4}{p{0.6\linewidth}}{
        \cellcolor[rgb]{0.9,0.9,0.9}{
            \makecell[{{p{\linewidth}}}]{
                \texttt{\tiny{[GM$\rangle$P2]}}
                \texttt{أنا مدير لعبتك، وأنت تلعب لعبة تعاونية مع الشبكة التالية كلوحة لعب:} \\
\\ 
\texttt{```} \\
\texttt{ 1234567} \\
\texttt{╔═╤══╤══╗} \\
\texttt{║◌│O◌│◌◌║ 1} \\
\texttt{║◌└┬─┴──╢ 2} \\
\texttt{║◌P│◌D◌◌║ 3} \\
\texttt{╟─┬┘◌◌┌─╢ 4} \\
\texttt{║◌│◌◌◌│A║ 5} \\
\texttt{╟─┼─┬─┤S║ 6} \\
\texttt{║M│◌│◌│T║ 7} \\
\texttt{╚═╧═╧═╧═╝} \\
\\ 
\texttt{```} \\
\\ 
\texttt{* تعرض الحافة العلوية إحداثيات x التي تتزايد نحو اليمين، والحافة اليمنى إحداثيات y التي تتزايد نحو الأسفل.} \\
\texttt{* الكائنات التالية موضوعة عشوائياً على شبكتك: 'P', 'O', 'T', 'S', 'D', 'A', 'M'.} \\
\\ 
\\ 
\texttt{يرى اللاعب الآخر نسخة مختلفة من لوحة اللعب، حيث توضع الكائنات في مواقع عشوائية مختلفة. لا يمكنك رؤية لوحة اللاعب الآخر، ولا يمكنه رؤية لوحتك.} \\
\\ 
\\ 
\texttt{**الهدف:**} \\
\\ 
\texttt{يجب على كلا اللاعبين تحريك الكائنات على خلفية كل منهما بحيث تنتهي الكائنات المتطابقة عند الإحداثيات نفسها. عليك التواصل مع اللاعب الآخر للاتفاق على إعداد هدف مشترك.} \\
\\ 
\\ 
\texttt{**القواعد:**} \\
\\ 
\texttt{* في كل دور، يمكنك إرسال أمر واحد فقط من الأمرين التاليين:} \\
\texttt{ 1. `قل: <رسالة>`: لإرسال رسالة (كل شيء حتى فاصل السطر التالي) إلى اللاعب الآخر. سأقوم بتمريرها إلى شريكك.} \\
\texttt{2. `حرّك: <كائن>, (<X>, <Y>)`: لتحريك كائن إلى موضع جديد، حيث `<X>` هو العمود و `<Y>` هو الصف. سأخبرك بما إذا كانت حركتك ناجحة أم لا.} \\
\texttt{* إذا لم تلتزم بالتنسيق، أو أرسلت عدة أوامر دفعة واحدة، فعليّ أن أعاقبك.} \\
\texttt{* إذا تجاوز كلا اللاعبين أكثر من 16 عقوبة، فستخسران اللعبة معاً.} \\
\texttt{* من الضروري أن تتواصل مع اللاعب الآخر بشأن حالة هدفك! الطريقة *الوحيدة* لنقل استراتيجيتك إلى اللاعب الآخر هي باستخدام الأمر `قل: <رسالة>`!} \\
\\ 
\\ 
\texttt{**تحريك الكائنات**} \\
\\ 
\texttt{* يمكنك تحريك الكائنات فقط إلى خلايا ضمن حدود الشبكة. يجب أن تكون الخلية المستهدفة فارغة، أي يجب أن تحتوي فقط على الرمز '◌'.} \\
\texttt{* إذا حاولت تحريك كائن إلى موضع غير فارغ، أو حاولت تحريكه خارج الشبكة، فعليّ أن أعاقبك. ستحصل على محاولة أخرى.} \\
\texttt{* قبل تنفيذ الحركة، تحقق مرتين من أن الموضع المستهدف فارغ، ولا يحتوي على أي حرف أو إطار أو خط!} \\
\\ 
\\ 
\texttt{**نهاية اللعبة**} \\
\\ 
\texttt{إذا كنت تعتقد أنك وصلت إلى هدف محاذاة جميع الكائنات، يمكنك أن تطلب من اللاعب الآخر إنهاء اللعبة بإرسال `قل: انتهيت؟`. إذا طلب منك اللاعب الآخر إنهاء اللعبة، ورددت بـ `قل: انتهيت!`، ستنتهي اللعبة.} \\
\\ 
\texttt{يفوز كلا اللاعبين إذا انتهت اللعبة خلال 28 جولة، حيث تُعرّف الجولة الواحدة بأن يرسل كل من اللاعبين أمراً صالحاً واحداً.} \\
\\ 
\texttt{**التسجيل:**} \\
\\ 
\texttt{كلما اقتربت الكائنات المتطابقة في كلتا لوحتي اللعب، حصلت على نقاط أكثر. العقوبات تقلل نقاطك. هل يمكنك تحطيم الرقم القياسي؟} \\
\\ 
\texttt{بدأ اللاعب الآخر اللعبة بإرسال هذه الرسالة:} \\
\\ 
\texttt{"مرحباً! لنبدأ بالتعاون. سأشارك مواقع الكائنات في لوحتي لنحدد هدفاً مشتركاً. T:(3,1), O:(3,3), P:(4,4), A:(2,5), D:(3,5), S:(2,7), M:(7,6). ما هي مواقع الكائنات لديك؟"} \\
\\ 
\texttt{ما هو أمرك الأول؟} \\
            }
        }
    }
    & \\ \\

    \theutterance \stepcounter{utterance}  
    & & & \multicolumn{4}{p{0.6\linewidth}}{
        \cellcolor[rgb]{0.9,0.9,0.9}{
            \makecell[{{p{\linewidth}}}]{
                \texttt{\tiny{[GM$\langle$P2]}}
                \texttt{قل: مرحباً! أنا مستعد للتعاون. إليك مواقع الكائنات في لوحتي:} \\
\texttt{P:(2,3), O:(3,1), T:(7,7), S:(7,6), D:(4,3), A:(7,5), M:(1,7).} \\
\texttt{لنحاول الاتفاق على هدف مشترك بناءً على هذه المواقع.} \\
            }
        }
    }
     \\ \\

    \theutterance \stepcounter{utterance}  
    & & & \multicolumn{2}{p{0.3\linewidth}}{
        \cellcolor[rgb]{0.9,0.9,0.9}{
            \makecell[{{p{\linewidth}}}]{
                \texttt{\tiny{[GM$|$GM]}}
                \texttt{قل: مرحباً! أنا مستعد للتعاون. إليك مواقع الكائنات في لوحتي:} \\
\texttt{P:(2,3), O:(3,1), T:(7,7), S:(7,6), D:(4,3), A:(7,5), M:(1,7).} \\
\texttt{لنحاول الاتفاق على هدف مشترك بناءً على هذه المواقع.} \\
            }
        }
    }
    & & \\ \\

    \theutterance \stepcounter{utterance}  
    & & & \multicolumn{2}{p{0.3\linewidth}}{
        \cellcolor[rgb]{0.9,0.9,0.9}{
            \makecell[{{p{\linewidth}}}]{
                \texttt{\tiny{[GM$|$GM]}}
                \texttt{Invalid format: tail is not empty: '} \\
\texttt{P:(2,3), O:(3,1), T:(7,7), S:(7,6), D:(4,3), A:(7,5), M:(1,7).} \\
\texttt{لنحاول الاتفاق على هدف مشترك بناءً على هذه المواقع.} \\
            }
        }
    }
    & & \\ \\

    \theutterance \stepcounter{utterance}  
    & & \multicolumn{4}{p{0.6\linewidth}}{
        \cellcolor[rgb]{0.9,0.9,0.9}{
            \makecell[{{p{\linewidth}}}]{
                \texttt{\tiny{[GM$\rangle$P2]}}
                \texttt{عقوبة: يجب ألا تحتوي رسالتك على أي شيء بعد الأمر!} \\
\\ 
\texttt{تأكد من أن ردك يحتوي فقط على `قل: <رسالة>` أو `حرّك: <كائن>, (<X>, <Y>)`، ولا شيء غير ذلك!} \\
\\ 
\texttt{لقد جمعتما معاً 3 من أصل 16 عقوبة.} \\
\texttt{يرجى المحاولة مرة أخرى!} \\
            }
        }
    }
    & \\ \\

    \theutterance \stepcounter{utterance}  
    & & & \multicolumn{4}{p{0.6\linewidth}}{
        \cellcolor[rgb]{0.9,0.9,0.9}{
            \makecell[{{p{\linewidth}}}]{
                \texttt{\tiny{[GM$\langle$P2]}}
                \texttt{قل: مرحباً! أنا مستعد للتعاون. إليك مواقع الكائنات في لوحتي: P:(2,3), O:(3,1), T:(7,7), S:(7,6), D:(4,3), A:(7,5), M:(1,7).} \\
            }
        }
    }
     \\ \\

    \theutterance \stepcounter{utterance}  
    & & & \multicolumn{2}{p{0.3\linewidth}}{
        \cellcolor[rgb]{0.9,0.9,0.9}{
            \makecell[{{p{\linewidth}}}]{
                \texttt{\tiny{[GM$|$GM]}}
                \texttt{قل: مرحباً! أنا مستعد للتعاون. إليك مواقع الكائنات في لوحتي: P:(2,3), O:(3,1), T:(7,7), S:(7,6), D:(4,3), A:(7,5), M:(1,7).} \\
            }
        }
    }
    & & \\ \\

    \theutterance \stepcounter{utterance}  
    & & \multicolumn{4}{p{0.6\linewidth}}{
        \cellcolor[rgb]{0.9,0.9,0.9}{
            \makecell[{{p{\linewidth}}}]{
                \texttt{\tiny{[P1$\langle$GM]}}
                \texttt{تم تمرير رسالتك إلى اللاعب الآخر.} \\
\\ 
\\ 
\texttt{أنت تلعب حالياً الجولة 1 من أصل 28 كحد أقصى.} \\
\texttt{لقد جمعتما معاً 3 من أصل 16 عقوبة.} \\
\texttt{أرسل اللاعب الآخر هذه الرسالة:} \\
\\ 
\\ 
\texttt{"مرحباً! أنا مستعد للتعاون. إليك مواقع الكائنات في لوحتي: P:(2,3), O:(3,1), T:(7,7), S:(7,6), D:(4,3), A:(7,5), M:(1,7)."} \\
\\ 
\\ 
\texttt{ما هو أمرك التالي؟} \\
            }
        }
    }
    & \\ \\

    \theutterance \stepcounter{utterance}  
    & \multicolumn{4}{p{0.6\linewidth}}{
        \cellcolor[rgb]{0.9,0.9,0.9}{
            \makecell[{{p{\linewidth}}}]{
                \texttt{\tiny{[P1$\rangle$GM]}}
                \texttt{قل: شكراً لك! لنحاول الاتفاق على هدف مشترك. ما رأيك أن نحاول نقل جميع الكائنات إلى مواقع محددة؟ لنقترح هذه الإحداثيات كهدف: T:(3,1), O:(3,3), P:(4,4), A:(2,5), D:(3,5), S:(2,7), M:(7,6). هل هذه المواقع متاحة في لوحتك؟} \\
            }
        }
    }
    & & \\ \\

    \theutterance \stepcounter{utterance}  
    & & & \multicolumn{2}{p{0.3\linewidth}}{
        \cellcolor[rgb]{0.9,0.9,0.9}{
            \makecell[{{p{\linewidth}}}]{
                \texttt{\tiny{[GM$|$GM]}}
                \texttt{قل: شكراً لك! لنحاول الاتفاق على هدف مشترك. ما رأيك أن نحاول نقل جميع الكائنات إلى مواقع محددة؟ لنقترح هذه الإحداثيات كهدف: T:(3,1), O:(3,3), P:(4,4), A:(2,5), D:(3,5), S:(2,7), M:(7,6). هل هذه المواقع متاحة في لوحتك؟} \\
            }
        }
    }
    & & \\ \\

    \theutterance \stepcounter{utterance}  
    & & \multicolumn{4}{p{0.6\linewidth}}{
        \cellcolor[rgb]{0.9,0.9,0.9}{
            \makecell[{{p{\linewidth}}}]{
                \texttt{\tiny{[GM$\rangle$P2]}}
                \texttt{تم تمرير رسالتك إلى اللاعب الآخر.} \\
\\ 
\\ 
\texttt{أنت تلعب حالياً الجولة 2 من أصل 28 كحد أقصى.} \\
\texttt{لقد جمعتما معاً 3 من أصل 16 عقوبة.} \\
\texttt{أرسل اللاعب الآخر هذه الرسالة:} \\
\\ 
\\ 
\texttt{"شكراً لك! لنحاول الاتفاق على هدف مشترك. ما رأيك أن نحاول نقل جميع الكائنات إلى مواقع محددة؟ لنقترح هذه الإحداثيات كهدف: T:(3,1), O:(3,3), P:(4,4), A:(2,5), D:(3,5), S:(2,7), M:(7,6). هل هذه المواقع متاحة في لوحتك؟"} \\
\\ 
\\ 
\texttt{ما هو أمرك التالي؟} \\
            }
        }
    }
    & \\ \\

    \theutterance \stepcounter{utterance}  
    & & & \multicolumn{4}{p{0.6\linewidth}}{
        \cellcolor[rgb]{0.9,0.9,0.9}{
            \makecell[{{p{\linewidth}}}]{
                \texttt{\tiny{[GM$\langle$P2]}}
                \texttt{قل: المواقع التي اقترحتها هي مواقع كائناتك، دعنا نتحقق من توفرها في لوحتي. المواقع المتاحة (الفارغة) هي: (3,1) غير متاح لأن O موجود هناك، (3,3) متاح، (4,4) متاح، (2,5) متاح، (3,5) متاح، (2,7) متاح، (7,6) متاح. هل يمكننا تعديل الهدف ليكون: T:(7,7), O:(3,1), P:(2,3), A:(7,5), D:(4,3), S:(7,6), M:(1,7)؟ (هذه هي مواقعي الحالية).} \\
            }
        }
    }
     \\ \\

    \theutterance \stepcounter{utterance}  
    & & & \multicolumn{2}{p{0.3\linewidth}}{
        \cellcolor[rgb]{0.9,0.9,0.9}{
            \makecell[{{p{\linewidth}}}]{
                \texttt{\tiny{[GM$|$GM]}}
                \texttt{قل: المواقع التي اقترحتها هي مواقع كائناتك، دعنا نتحقق من توفرها في لوحتي. المواقع المتاحة (الفارغة) هي: (3,1) غير متاح لأن O موجود هناك، (3,3) متاح، (4,4) متاح، (2,5) متاح، (3,5) متاح، (2,7) متاح، (7,6) متاح. هل يمكننا تعديل الهدف ليكون: T:(7,7), O:(3,1), P:(2,3), A:(7,5), D:(4,3), S:(7,6), M:(1,7)؟ (هذه هي مواقعي الحالية).} \\
            }
        }
    }
    & & \\ \\

    \theutterance \stepcounter{utterance}  
    & & \multicolumn{4}{p{0.6\linewidth}}{
        \cellcolor[rgb]{0.9,0.9,0.9}{
            \makecell[{{p{\linewidth}}}]{
                \texttt{\tiny{[P1$\langle$GM]}}
                \texttt{تم تمرير رسالتك إلى اللاعب الآخر.} \\
\\ 
\\ 
\texttt{أنت تلعب حالياً الجولة 2 من أصل 28 كحد أقصى.} \\
\texttt{لقد جمعتما معاً 3 من أصل 16 عقوبة.} \\
\texttt{أرسل اللاعب الآخر هذه الرسالة:} \\
\\ 
\\ 
\texttt{"المواقع التي اقترحتها هي مواقع كائناتك، دعنا نتحقق من توفرها في لوحتي. المواقع المتاحة (الفارغة) هي: (3,1) غير متاح لأن O موجود هناك، (3,3) متاح، (4,4) متاح، (2,5) متاح، (3,5) متاح، (2,7) متاح، (7,6) متاح. هل يمكننا تعديل الهدف ليكون: T:(7,7), O:(3,1), P:(2,3), A:(7,5), D:(4,3), S:(7,6), M:(1,7)؟ (هذه هي مواقعي الحالية)."} \\
\\ 
\\ 
\texttt{ما هو أمرك التالي؟} \\
            }
        }
    }
    & \\ \\

    \theutterance \stepcounter{utterance}  
    & \multicolumn{4}{p{0.6\linewidth}}{
        \cellcolor[rgb]{0.9,0.9,0.9}{
            \makecell[{{p{\linewidth}}}]{
                \texttt{\tiny{[P1$\rangle$GM]}}
                \texttt{قل: اقتراح جيد! لنعتمد هدفك (مواقع كائناتك الحالية) كهدف مشترك لنا. سأبدأ بتحريك كائناتي لتطابق إحداثياتك. سأبدأ بـ T إلى (7, 7).} \\
            }
        }
    }
    & & \\ \\

    \theutterance \stepcounter{utterance}  
    & & & \multicolumn{2}{p{0.3\linewidth}}{
        \cellcolor[rgb]{0.9,0.9,0.9}{
            \makecell[{{p{\linewidth}}}]{
                \texttt{\tiny{[GM$|$GM]}}
                \texttt{قل: اقتراح جيد! لنعتمد هدفك (مواقع كائناتك الحالية) كهدف مشترك لنا. سأبدأ بتحريك كائناتي لتطابق إحداثياتك. سأبدأ بـ T إلى (7, 7).} \\
            }
        }
    }
    & & \\ \\

    \theutterance \stepcounter{utterance}  
    & & \multicolumn{4}{p{0.6\linewidth}}{
        \cellcolor[rgb]{0.9,0.9,0.9}{
            \makecell[{{p{\linewidth}}}]{
                \texttt{\tiny{[GM$\rangle$P2]}}
                \texttt{تم تمرير رسالتك إلى اللاعب الآخر.} \\
\\ 
\\ 
\texttt{أنت تلعب حالياً الجولة 3 من أصل 28 كحد أقصى.} \\
\texttt{لقد جمعتما معاً 3 من أصل 16 عقوبة.} \\
\texttt{أرسل اللاعب الآخر هذه الرسالة:} \\
\\ 
\\ 
\texttt{"اقتراح جيد! لنعتمد هدفك (مواقع كائناتك الحالية) كهدف مشترك لنا. سأبدأ بتحريك كائناتي لتطابق إحداثياتك. سأبدأ بـ T إلى (7, 7)."} \\
\\ 
\\ 
\texttt{ما هو أمرك التالي؟} \\
            }
        }
    }
    & \\ \\

    \theutterance \stepcounter{utterance}  
    & & & \multicolumn{4}{p{0.6\linewidth}}{
        \cellcolor[rgb]{0.9,0.9,0.9}{
            \makecell[{{p{\linewidth}}}]{
                \texttt{\tiny{[GM$\langle$P2]}}
                \texttt{قل: ممتاز! سأبدأ أنا أيضاً بتحريك كائناتي لتطابق مواقعها الحالية (بما أنها هي الهدف بالفعل، سأبدأ بتحريك الكائنات التي قد تحتاج لتعديل في لوحتك أو التأكد من ثباتها). سأبدأ بتحريك O إلى (3,1) إذا لم يكن هناك عائق. (ملاحظة: بما أن O موجود بالفعل في (3,1)، سأنتظر حركتك القادمة أو أبدأ بتحريك كائن آخر إذا لزم الأمر). لنبدأ بـ O: (3,1).} \\
            }
        }
    }
     \\ \\

    \theutterance \stepcounter{utterance}  
    & & & \multicolumn{2}{p{0.3\linewidth}}{
        \cellcolor[rgb]{0.9,0.9,0.9}{
            \makecell[{{p{\linewidth}}}]{
                \texttt{\tiny{[GM$|$GM]}}
                \texttt{قل: ممتاز! سأبدأ أنا أيضاً بتحريك كائناتي لتطابق مواقعها الحالية (بما أنها هي الهدف بالفعل، سأبدأ بتحريك الكائنات التي قد تحتاج لتعديل في لوحتك أو التأكد من ثباتها). سأبدأ بتحريك O إلى (3,1) إذا لم يكن هناك عائق. (ملاحظة: بما أن O موجود بالفعل في (3,1)، سأنتظر حركتك القادمة أو أبدأ بتحريك كائن آخر إذا لزم الأمر). لنبدأ بـ O: (3,1).} \\
            }
        }
    }
    & & \\ \\

    \theutterance \stepcounter{utterance}  
    & & \multicolumn{4}{p{0.6\linewidth}}{
        \cellcolor[rgb]{0.9,0.9,0.9}{
            \makecell[{{p{\linewidth}}}]{
                \texttt{\tiny{[P1$\langle$GM]}}
                \texttt{تم تمرير رسالتك إلى اللاعب الآخر.} \\
\\ 
\\ 
\texttt{أنت تلعب حالياً الجولة 3 من أصل 28 كحد أقصى.} \\
\texttt{لقد جمعتما معاً 3 من أصل 16 عقوبة.} \\
\texttt{أرسل اللاعب الآخر هذه الرسالة:} \\
\\ 
\\ 
\texttt{"ممتاز! سأبدأ أنا أيضاً بتحريك كائناتي لتطابق مواقعها الحالية (بما أنها هي الهدف بالفعل، سأبدأ بتحريك الكائنات التي قد تحتاج لتعديل في لوحتك أو التأكد من ثباتها). سأبدأ بتحريك O إلى (3,1) إذا لم يكن هناك عائق. (ملاحظة: بما أن O موجود بالفعل في (3,1)، سأنتظر حركتك القادمة أو أبدأ بتحريك كائن آخر إذا لزم الأمر). لنبدأ بـ O: (3,1)."} \\
\\ 
\\ 
\texttt{ما هو أمرك التالي؟} \\
            }
        }
    }
    & \\ \\

    \theutterance \stepcounter{utterance}  
    & \multicolumn{4}{p{0.6\linewidth}}{
        \cellcolor[rgb]{0.9,0.9,0.9}{
            \makecell[{{p{\linewidth}}}]{
                \texttt{\tiny{[P1$\rangle$GM]}}
                \texttt{حرّك: T, (7, 7)} \\
            }
        }
    }
    & & \\ \\

    \theutterance \stepcounter{utterance}  
    & & & \multicolumn{2}{p{0.3\linewidth}}{
        \cellcolor[rgb]{0.9,0.9,0.9}{
            \makecell[{{p{\linewidth}}}]{
                \texttt{\tiny{[GM$|$GM]}}
                \texttt{حرّك: T, (7, 7)} \\
            }
        }
    }
    & & \\ \\

    \theutterance \stepcounter{utterance}  
    & & & \multicolumn{2}{p{0.3\linewidth}}{
        \cellcolor[rgb]{0.9,0.9,0.9}{
            \makecell[{{p{\linewidth}}}]{
                \texttt{\tiny{[GM$|$GM]}}
                \texttt{T moved to (7, 7)} \\
            }
        }
    }
    & & \\ \\

    \theutterance \stepcounter{utterance}  
    & & \multicolumn{4}{p{0.6\linewidth}}{
        \cellcolor[rgb]{0.9,0.9,0.9}{
            \makecell[{{p{\linewidth}}}]{
                \texttt{\tiny{[GM$\rangle$P2]}}
                \texttt{تم تمرير رسالتك إلى اللاعب الآخر.} \\
\\ 
\\ 
\texttt{أنت تلعب حالياً الجولة 4 من أصل 28 كحد أقصى.} \\
\texttt{لقد جمعتما معاً 3 من أصل 16 عقوبة.} \\
\texttt{حرّك اللاعب الآخر كائناً على شبكته.} \\
\\ 
\\ 
\texttt{ما هو أمرك التالي؟} \\
            }
        }
    }
    & \\ \\

    \theutterance \stepcounter{utterance}  
    & & & \multicolumn{4}{p{0.6\linewidth}}{
        \cellcolor[rgb]{0.9,0.9,0.9}{
            \makecell[{{p{\linewidth}}}]{
                \texttt{\tiny{[GM$\langle$P2]}}
                \texttt{حرّك: P, (2, 3)} \\
            }
        }
    }
     \\ \\

    \theutterance \stepcounter{utterance}  
    & & & \multicolumn{2}{p{0.3\linewidth}}{
        \cellcolor[rgb]{0.9,0.9,0.9}{
            \makecell[{{p{\linewidth}}}]{
                \texttt{\tiny{[GM$|$GM]}}
                \texttt{حرّك: P, (2, 3)} \\
            }
        }
    }
    & & \\ \\

    \theutterance \stepcounter{utterance}  
    & & & \multicolumn{2}{p{0.3\linewidth}}{
        \cellcolor[rgb]{0.9,0.9,0.9}{
            \makecell[{{p{\linewidth}}}]{
                \texttt{\tiny{[GM$|$GM]}}
                \texttt{عقوبة: (2,3) ليست فارغة، بل تحتوي على 'P'.} \\
\\ 
\texttt{لقد جمعتما معاً 4 من أصل 16 عقوبة.} \\
\texttt{يرجى المحاولة مرة أخرى!} \\
            }
        }
    }
    & & \\ \\

    \theutterance \stepcounter{utterance}  
    & & \multicolumn{4}{p{0.6\linewidth}}{
        \cellcolor[rgb]{0.9,0.9,0.9}{
            \makecell[{{p{\linewidth}}}]{
                \texttt{\tiny{[GM$\rangle$P2]}}
                \texttt{عقوبة: (2,3) ليست فارغة، بل تحتوي على 'P'.} \\
\\ 
\texttt{لقد جمعتما معاً 4 من أصل 16 عقوبة.} \\
\texttt{يرجى المحاولة مرة أخرى!} \\
            }
        }
    }
    & \\ \\

    \theutterance \stepcounter{utterance}  
    & & & \multicolumn{4}{p{0.6\linewidth}}{
        \cellcolor[rgb]{0.9,0.9,0.9}{
            \makecell[{{p{\linewidth}}}]{
                \texttt{\tiny{[GM$\langle$P2]}}
                \texttt{حرّك: O, (3, 1)} \\
            }
        }
    }
     \\ \\

    \theutterance \stepcounter{utterance}  
    & & & \multicolumn{2}{p{0.3\linewidth}}{
        \cellcolor[rgb]{0.9,0.9,0.9}{
            \makecell[{{p{\linewidth}}}]{
                \texttt{\tiny{[GM$|$GM]}}
                \texttt{حرّك: O, (3, 1)} \\
            }
        }
    }
    & & \\ \\

    \theutterance \stepcounter{utterance}  
    & & & \multicolumn{2}{p{0.3\linewidth}}{
        \cellcolor[rgb]{0.9,0.9,0.9}{
            \makecell[{{p{\linewidth}}}]{
                \texttt{\tiny{[GM$|$GM]}}
                \texttt{عقوبة: (3,1) ليست فارغة، بل تحتوي على 'O'.} \\
\\ 
\texttt{لقد جمعتما معاً 5 من أصل 16 عقوبة.} \\
\texttt{يرجى المحاولة مرة أخرى!} \\
            }
        }
    }
    & & \\ \\

    \theutterance \stepcounter{utterance}  
    & & \multicolumn{4}{p{0.6\linewidth}}{
        \cellcolor[rgb]{0.9,0.9,0.9}{
            \makecell[{{p{\linewidth}}}]{
                \texttt{\tiny{[GM$\rangle$P2]}}
                \texttt{عقوبة: (3,1) ليست فارغة، بل تحتوي على 'O'.} \\
\\ 
\texttt{لقد جمعتما معاً 5 من أصل 16 عقوبة.} \\
\texttt{يرجى المحاولة مرة أخرى!} \\
            }
        }
    }
    & \\ \\

    \theutterance \stepcounter{utterance}  
    & & & \multicolumn{4}{p{0.6\linewidth}}{
        \cellcolor[rgb]{0.9,0.9,0.9}{
            \makecell[{{p{\linewidth}}}]{
                \texttt{\tiny{[GM$\langle$P2]}}
                \texttt{حرّك: D, (4, 3)} \\
            }
        }
    }
     \\ \\

    \theutterance \stepcounter{utterance}  
    & & & \multicolumn{2}{p{0.3\linewidth}}{
        \cellcolor[rgb]{0.9,0.9,0.9}{
            \makecell[{{p{\linewidth}}}]{
                \texttt{\tiny{[GM$|$GM]}}
                \texttt{حرّك: D, (4, 3)} \\
            }
        }
    }
    & & \\ \\

    \theutterance \stepcounter{utterance}  
    & & & \multicolumn{2}{p{0.3\linewidth}}{
        \cellcolor[rgb]{0.9,0.9,0.9}{
            \makecell[{{p{\linewidth}}}]{
                \texttt{\tiny{[GM$|$GM]}}
                \texttt{D moved to (4, 3)} \\
            }
        }
    }
    & & \\ \\

    \theutterance \stepcounter{utterance}  
    & & \multicolumn{4}{p{0.6\linewidth}}{
        \cellcolor[rgb]{0.9,0.9,0.9}{
            \makecell[{{p{\linewidth}}}]{
                \texttt{\tiny{[P1$\langle$GM]}}
                \texttt{تم تحريك 'T' إلى (7,7) بنجاح. تبدو شبكتك المحدّثة هكذا:} \\
\texttt{```} \\
\texttt{ 1234567} \\
\texttt{╔═╤══╤══╗} \\
\texttt{║◌│◌◌│◌◌║ 1} \\
\texttt{║◌└┬─┴──╢ 2} \\
\texttt{║◌◌│O◌◌◌║ 3} \\
\texttt{╟─┬┘P◌┌─╢ 4} \\
\texttt{║◌│AD◌│◌║ 5} \\
\texttt{╟─┼─┬─┤M║ 6} \\
\texttt{║◌│S│◌│T║ 7} \\
\texttt{╚═╧═╧═╧═╝} \\
\texttt{```} \\
\texttt{أنت تلعب حالياً الجولة 4 من أصل 28 كحد أقصى.} \\
\texttt{لقد جمعتما معاً 5 من أصل 16 عقوبة.} \\
\texttt{حرّك اللاعب الآخر كائناً على شبكته.} \\
\\ 
\\ 
\texttt{ما هو أمرك التالي؟} \\
            }
        }
    }
    & \\ \\

    \theutterance \stepcounter{utterance}  
    & \multicolumn{4}{p{0.6\linewidth}}{
        \cellcolor[rgb]{0.9,0.9,0.9}{
            \makecell[{{p{\linewidth}}}]{
                \texttt{\tiny{[P1$\rangle$GM]}}
                \texttt{حرّك: O, (3, 1)} \\
            }
        }
    }
    & & \\ \\

    \theutterance \stepcounter{utterance}  
    & & & \multicolumn{2}{p{0.3\linewidth}}{
        \cellcolor[rgb]{0.9,0.9,0.9}{
            \makecell[{{p{\linewidth}}}]{
                \texttt{\tiny{[GM$|$GM]}}
                \texttt{حرّك: O, (3, 1)} \\
            }
        }
    }
    & & \\ \\

    \theutterance \stepcounter{utterance}  
    & & & \multicolumn{2}{p{0.3\linewidth}}{
        \cellcolor[rgb]{0.9,0.9,0.9}{
            \makecell[{{p{\linewidth}}}]{
                \texttt{\tiny{[GM$|$GM]}}
                \texttt{O moved to (3, 1)} \\
            }
        }
    }
    & & \\ \\

    \theutterance \stepcounter{utterance}  
    & & \multicolumn{4}{p{0.6\linewidth}}{
        \cellcolor[rgb]{0.9,0.9,0.9}{
            \makecell[{{p{\linewidth}}}]{
                \texttt{\tiny{[GM$\rangle$P2]}}
                \texttt{تم تحريك 'D' إلى (4,3) بنجاح. تبدو شبكتك المحدّثة هكذا:} \\
\texttt{```} \\
\texttt{ 1234567} \\
\texttt{╔═╤══╤══╗} \\
\texttt{║◌│O◌│◌◌║ 1} \\
\texttt{║◌└┬─┴──╢ 2} \\
\texttt{║◌P│D◌◌◌║ 3} \\
\texttt{╟─┬┘◌◌┌─╢ 4} \\
\texttt{║◌│◌◌◌│A║ 5} \\
\texttt{╟─┼─┬─┤S║ 6} \\
\texttt{║M│◌│◌│T║ 7} \\
\texttt{╚═╧═╧═╧═╝} \\
\texttt{```} \\
\texttt{أنت تلعب حالياً الجولة 5 من أصل 28 كحد أقصى.} \\
\texttt{لقد جمعتما معاً 5 من أصل 16 عقوبة.} \\
\texttt{حرّك اللاعب الآخر كائناً على شبكته.} \\
\\ 
\\ 
\texttt{ما هو أمرك التالي؟} \\
            }
        }
    }
    & \\ \\

    \theutterance \stepcounter{utterance}  
    & & & \multicolumn{4}{p{0.6\linewidth}}{
        \cellcolor[rgb]{0.9,0.9,0.9}{
            \makecell[{{p{\linewidth}}}]{
                \texttt{\tiny{[GM$\langle$P2]}}
                \texttt{حرّك: A, (7, 5)} \\
            }
        }
    }
     \\ \\

    \theutterance \stepcounter{utterance}  
    & & & \multicolumn{2}{p{0.3\linewidth}}{
        \cellcolor[rgb]{0.9,0.9,0.9}{
            \makecell[{{p{\linewidth}}}]{
                \texttt{\tiny{[GM$|$GM]}}
                \texttt{حرّك: A, (7, 5)} \\
            }
        }
    }
    & & \\ \\

    \theutterance \stepcounter{utterance}  
    & & & \multicolumn{2}{p{0.3\linewidth}}{
        \cellcolor[rgb]{0.9,0.9,0.9}{
            \makecell[{{p{\linewidth}}}]{
                \texttt{\tiny{[GM$|$GM]}}
                \texttt{عقوبة: (7,5) ليست فارغة، بل تحتوي على 'A'.} \\
\\ 
\texttt{لقد جمعتما معاً 6 من أصل 16 عقوبة.} \\
\texttt{يرجى المحاولة مرة أخرى!} \\
            }
        }
    }
    & & \\ \\

    \theutterance \stepcounter{utterance}  
    & & \multicolumn{4}{p{0.6\linewidth}}{
        \cellcolor[rgb]{0.9,0.9,0.9}{
            \makecell[{{p{\linewidth}}}]{
                \texttt{\tiny{[GM$\rangle$P2]}}
                \texttt{عقوبة: (7,5) ليست فارغة، بل تحتوي على 'A'.} \\
\\ 
\texttt{لقد جمعتما معاً 6 من أصل 16 عقوبة.} \\
\texttt{يرجى المحاولة مرة أخرى!} \\
            }
        }
    }
    & \\ \\

    \theutterance \stepcounter{utterance}  
    & & & \multicolumn{4}{p{0.6\linewidth}}{
        \cellcolor[rgb]{0.9,0.9,0.9}{
            \makecell[{{p{\linewidth}}}]{
                \texttt{\tiny{[GM$\langle$P2]}}
                \texttt{حرّك: M, (1, 7)} \\
            }
        }
    }
     \\ \\

    \theutterance \stepcounter{utterance}  
    & & & \multicolumn{2}{p{0.3\linewidth}}{
        \cellcolor[rgb]{0.9,0.9,0.9}{
            \makecell[{{p{\linewidth}}}]{
                \texttt{\tiny{[GM$|$GM]}}
                \texttt{حرّك: M, (1, 7)} \\
            }
        }
    }
    & & \\ \\

    \theutterance \stepcounter{utterance}  
    & & & \multicolumn{2}{p{0.3\linewidth}}{
        \cellcolor[rgb]{0.9,0.9,0.9}{
            \makecell[{{p{\linewidth}}}]{
                \texttt{\tiny{[GM$|$GM]}}
                \texttt{عقوبة: (1,7) ليست فارغة، بل تحتوي على 'M'.} \\
\\ 
\texttt{لقد جمعتما معاً 7 من أصل 16 عقوبة.} \\
\texttt{يرجى المحاولة مرة أخرى!} \\
            }
        }
    }
    & & \\ \\

    \theutterance \stepcounter{utterance}  
    & & \multicolumn{4}{p{0.6\linewidth}}{
        \cellcolor[rgb]{0.9,0.9,0.9}{
            \makecell[{{p{\linewidth}}}]{
                \texttt{\tiny{[GM$\rangle$P2]}}
                \texttt{عقوبة: (1,7) ليست فارغة، بل تحتوي على 'M'.} \\
\\ 
\texttt{لقد جمعتما معاً 7 من أصل 16 عقوبة.} \\
\texttt{يرجى المحاولة مرة أخرى!} \\
            }
        }
    }
    & \\ \\

    \theutterance \stepcounter{utterance}  
    & & & \multicolumn{4}{p{0.6\linewidth}}{
        \cellcolor[rgb]{0.9,0.9,0.9}{
            \makecell[{{p{\linewidth}}}]{
                \texttt{\tiny{[GM$\langle$P2]}}
                \texttt{حرّك: T, (7, 7)} \\
            }
        }
    }
     \\ \\

    \theutterance \stepcounter{utterance}  
    & & & \multicolumn{2}{p{0.3\linewidth}}{
        \cellcolor[rgb]{0.9,0.9,0.9}{
            \makecell[{{p{\linewidth}}}]{
                \texttt{\tiny{[GM$|$GM]}}
                \texttt{حرّك: T, (7, 7)} \\
            }
        }
    }
    & & \\ \\

    \theutterance \stepcounter{utterance}  
    & & & \multicolumn{2}{p{0.3\linewidth}}{
        \cellcolor[rgb]{0.9,0.9,0.9}{
            \makecell[{{p{\linewidth}}}]{
                \texttt{\tiny{[GM$|$GM]}}
                \texttt{عقوبة: (7,7) ليست فارغة، بل تحتوي على 'T'.} \\
\\ 
\texttt{لقد جمعتما معاً 8 من أصل 16 عقوبة.} \\
\texttt{يرجى المحاولة مرة أخرى!} \\
            }
        }
    }
    & & \\ \\

    \theutterance \stepcounter{utterance}  
    & & \multicolumn{4}{p{0.6\linewidth}}{
        \cellcolor[rgb]{0.9,0.9,0.9}{
            \makecell[{{p{\linewidth}}}]{
                \texttt{\tiny{[GM$\rangle$P2]}}
                \texttt{عقوبة: (7,7) ليست فارغة، بل تحتوي على 'T'.} \\
\\ 
\texttt{لقد جمعتما معاً 8 من أصل 16 عقوبة.} \\
\texttt{يرجى المحاولة مرة أخرى!} \\
            }
        }
    }
    & \\ \\

    \theutterance \stepcounter{utterance}  
    & & & \multicolumn{4}{p{0.6\linewidth}}{
        \cellcolor[rgb]{0.9,0.9,0.9}{
            \makecell[{{p{\linewidth}}}]{
                \texttt{\tiny{[GM$\langle$P2]}}
                \texttt{حرّك: S, (7, 6)} \\
            }
        }
    }
     \\ \\

    \theutterance \stepcounter{utterance}  
    & & & \multicolumn{2}{p{0.3\linewidth}}{
        \cellcolor[rgb]{0.9,0.9,0.9}{
            \makecell[{{p{\linewidth}}}]{
                \texttt{\tiny{[GM$|$GM]}}
                \texttt{حرّك: S, (7, 6)} \\
            }
        }
    }
    & & \\ \\

    \theutterance \stepcounter{utterance}  
    & & & \multicolumn{2}{p{0.3\linewidth}}{
        \cellcolor[rgb]{0.9,0.9,0.9}{
            \makecell[{{p{\linewidth}}}]{
                \texttt{\tiny{[GM$|$GM]}}
                \texttt{عقوبة: (7,6) ليست فارغة، بل تحتوي على 'S'.} \\
\\ 
\texttt{لقد جمعتما معاً 9 من أصل 16 عقوبة.} \\
\texttt{يرجى المحاولة مرة أخرى!} \\
            }
        }
    }
    & & \\ \\

    \theutterance \stepcounter{utterance}  
    & & \multicolumn{4}{p{0.6\linewidth}}{
        \cellcolor[rgb]{0.9,0.9,0.9}{
            \makecell[{{p{\linewidth}}}]{
                \texttt{\tiny{[GM$\rangle$P2]}}
                \texttt{عقوبة: (7,6) ليست فارغة، بل تحتوي على 'S'.} \\
\\ 
\texttt{لقد جمعتما معاً 9 من أصل 16 عقوبة.} \\
\texttt{يرجى المحاولة مرة أخرى!} \\
            }
        }
    }
    & \\ \\

    \theutterance \stepcounter{utterance}  
    & & & \multicolumn{4}{p{0.6\linewidth}}{
        \cellcolor[rgb]{0.9,0.9,0.9}{
            \makecell[{{p{\linewidth}}}]{
                \texttt{\tiny{[GM$\langle$P2]}}
                \texttt{حرّك: S, (6, 6)} \\
            }
        }
    }
     \\ \\

    \theutterance \stepcounter{utterance}  
    & & & \multicolumn{2}{p{0.3\linewidth}}{
        \cellcolor[rgb]{0.9,0.9,0.9}{
            \makecell[{{p{\linewidth}}}]{
                \texttt{\tiny{[GM$|$GM]}}
                \texttt{حرّك: S, (6, 6)} \\
            }
        }
    }
    & & \\ \\

    \theutterance \stepcounter{utterance}  
    & & & \multicolumn{2}{p{0.3\linewidth}}{
        \cellcolor[rgb]{0.9,0.9,0.9}{
            \makecell[{{p{\linewidth}}}]{
                \texttt{\tiny{[GM$|$GM]}}
                \texttt{عقوبة: (6,6) ليست فارغة، بل تحتوي على '┤'.} \\
\\ 
\texttt{لقد جمعتما معاً 10 من أصل 16 عقوبة.} \\
\texttt{يرجى المحاولة مرة أخرى!} \\
            }
        }
    }
    & & \\ \\

    \theutterance \stepcounter{utterance}  
    & & \multicolumn{4}{p{0.6\linewidth}}{
        \cellcolor[rgb]{0.9,0.9,0.9}{
            \makecell[{{p{\linewidth}}}]{
                \texttt{\tiny{[GM$\rangle$P2]}}
                \texttt{عقوبة: (6,6) ليست فارغة، بل تحتوي على '┤'.} \\
\\ 
\texttt{لقد جمعتما معاً 10 من أصل 16 عقوبة.} \\
\texttt{يرجى المحاولة مرة أخرى!} \\
            }
        }
    }
    & \\ \\

    \theutterance \stepcounter{utterance}  
    & & & \multicolumn{4}{p{0.6\linewidth}}{
        \cellcolor[rgb]{0.9,0.9,0.9}{
            \makecell[{{p{\linewidth}}}]{
                \texttt{\tiny{[GM$\langle$P2]}}
                \texttt{حرّك: S, (7, 6)} \\
\\ 
\texttt{(ملاحظة للمدير: يبدو أنني أخطأت في قراءة الإحداثيات في المحاولات السابقة. سأحاول تحريك S إلى مكان فارغ تماماً. بالنظر للشبكة، الخلية (7,6) تحتوي على 'S' بالفعل، والخلية (6,6) تحتوي على إطار. سأحاول تحريك S إلى (6,7) إذا كانت فارغة، أو سأحاول تحريك كائن آخر إلى مكان فارغ تماماً مثل (2,1) أو (4,1)).} \\
\\ 
\texttt{حرّك: S, (6, 7)} \\
            }
        }
    }
     \\ \\

    \theutterance \stepcounter{utterance}  
    & & & \multicolumn{2}{p{0.3\linewidth}}{
        \cellcolor[rgb]{0.9,0.9,0.9}{
            \makecell[{{p{\linewidth}}}]{
                \texttt{\tiny{[GM$|$GM]}}
                \texttt{حرّك: S, (7, 6)} \\
\\ 
\texttt{(ملاحظة للمدير: يبدو أنني أخطأت في قراءة الإحداثيات في المحاولات السابقة. سأحاول تحريك S إلى مكان فارغ تماماً. بالنظر للشبكة، الخلية (7,6) تحتوي على 'S' بالفعل، والخلية (6,6) تحتوي على إطار. سأحاول تحريك S إلى (6,7) إذا كانت فارغة، أو سأحاول تحريك كائن آخر إلى مكان فارغ تماماً مثل (2,1) أو (4,1)).} \\
\\ 
\texttt{حرّك: S, (6, 7)} \\
            }
        }
    }
    & & \\ \\

    \theutterance \stepcounter{utterance}  
    & & & \multicolumn{2}{p{0.3\linewidth}}{
        \cellcolor[rgb]{0.9,0.9,0.9}{
            \makecell[{{p{\linewidth}}}]{
                \texttt{\tiny{[GM$|$GM]}}
                \texttt{Invalid format: head is not empty: 'حرّك: S, (7, 6)} \\
\\ 
\texttt{(ملاحظة للمدير: يبدو أنني أخطأت في قراءة الإحداثيات في المحاولات السابقة. سأحاول تحريك S إلى مكان فارغ تماماً. بالنظر للشبكة، الخلية (7,6) تحتوي على 'S' بالفعل، والخلية (6,6) تحتوي على إطار. سأحاول تحريك S إلى (6,7) إذا كانت فارغة، أو سأحاول تحريك كائن آخر إلى مكان فارغ تماماً مثل (2,1) أو (4,1)).} \\
\\ 
\texttt{'} \\
            }
        }
    }
    & & \\ \\

    \theutterance \stepcounter{utterance}  
    & & \multicolumn{4}{p{0.6\linewidth}}{
        \cellcolor[rgb]{0.9,0.9,0.9}{
            \makecell[{{p{\linewidth}}}]{
                \texttt{\tiny{[GM$\rangle$P2]}}
                \texttt{عقوبة: يجب ألا تحتوي رسالتك على أي شيء قبل الأمر!} \\
\\ 
\texttt{تأكد من أن ردك يحتوي فقط على `قل: <رسالة>` أو `حرّك: <كائن>, (<X>, <Y>)`، ولا شيء غير ذلك!} \\
\\ 
\texttt{لقد جمعتما معاً 11 من أصل 16 عقوبة.} \\
\texttt{يرجى المحاولة مرة أخرى!} \\
            }
        }
    }
    & \\ \\

    \theutterance \stepcounter{utterance}  
    & & & \multicolumn{4}{p{0.6\linewidth}}{
        \cellcolor[rgb]{0.9,0.9,0.9}{
            \makecell[{{p{\linewidth}}}]{
                \texttt{\tiny{[GM$\langle$P2]}}
                \texttt{حرّك: S, (6, 7)} \\
            }
        }
    }
     \\ \\

    \theutterance \stepcounter{utterance}  
    & & & \multicolumn{2}{p{0.3\linewidth}}{
        \cellcolor[rgb]{0.9,0.9,0.9}{
            \makecell[{{p{\linewidth}}}]{
                \texttt{\tiny{[GM$|$GM]}}
                \texttt{حرّك: S, (6, 7)} \\
            }
        }
    }
    & & \\ \\

    \theutterance \stepcounter{utterance}  
    & & & \multicolumn{2}{p{0.3\linewidth}}{
        \cellcolor[rgb]{0.9,0.9,0.9}{
            \makecell[{{p{\linewidth}}}]{
                \texttt{\tiny{[GM$|$GM]}}
                \texttt{عقوبة: (6,7) ليست فارغة، بل تحتوي على '│'.} \\
\\ 
\texttt{لقد جمعتما معاً 12 من أصل 16 عقوبة.} \\
\texttt{يرجى المحاولة مرة أخرى!} \\
            }
        }
    }
    & & \\ \\

    \theutterance \stepcounter{utterance}  
    & & \multicolumn{4}{p{0.6\linewidth}}{
        \cellcolor[rgb]{0.9,0.9,0.9}{
            \makecell[{{p{\linewidth}}}]{
                \texttt{\tiny{[GM$\rangle$P2]}}
                \texttt{عقوبة: (6,7) ليست فارغة، بل تحتوي على '│'.} \\
\\ 
\texttt{لقد جمعتما معاً 12 من أصل 16 عقوبة.} \\
\texttt{يرجى المحاولة مرة أخرى!} \\
            }
        }
    }
    & \\ \\

    \theutterance \stepcounter{utterance}  
    & & & \multicolumn{4}{p{0.6\linewidth}}{
        \cellcolor[rgb]{0.9,0.9,0.9}{
            \makecell[{{p{\linewidth}}}]{
                \texttt{\tiny{[GM$\langle$P2]}}
                \texttt{حرّك: S, (5, 7)} \\
            }
        }
    }
     \\ \\

    \theutterance \stepcounter{utterance}  
    & & & \multicolumn{2}{p{0.3\linewidth}}{
        \cellcolor[rgb]{0.9,0.9,0.9}{
            \makecell[{{p{\linewidth}}}]{
                \texttt{\tiny{[GM$|$GM]}}
                \texttt{حرّك: S, (5, 7)} \\
            }
        }
    }
    & & \\ \\

    \theutterance \stepcounter{utterance}  
    & & & \multicolumn{2}{p{0.3\linewidth}}{
        \cellcolor[rgb]{0.9,0.9,0.9}{
            \makecell[{{p{\linewidth}}}]{
                \texttt{\tiny{[GM$|$GM]}}
                \texttt{S moved to (5, 7)} \\
            }
        }
    }
    & & \\ \\

    \theutterance \stepcounter{utterance}  
    & & \multicolumn{4}{p{0.6\linewidth}}{
        \cellcolor[rgb]{0.9,0.9,0.9}{
            \makecell[{{p{\linewidth}}}]{
                \texttt{\tiny{[P1$\langle$GM]}}
                \texttt{تم تحريك 'O' إلى (3,1) بنجاح. تبدو شبكتك المحدّثة هكذا:} \\
\texttt{```} \\
\texttt{ 1234567} \\
\texttt{╔═╤══╤══╗} \\
\texttt{║◌│O◌│◌◌║ 1} \\
\texttt{║◌└┬─┴──╢ 2} \\
\texttt{║◌◌│◌◌◌◌║ 3} \\
\texttt{╟─┬┘P◌┌─╢ 4} \\
\texttt{║◌│AD◌│◌║ 5} \\
\texttt{╟─┼─┬─┤M║ 6} \\
\texttt{║◌│S│◌│T║ 7} \\
\texttt{╚═╧═╧═╧═╝} \\
\texttt{```} \\
\texttt{أنت تلعب حالياً الجولة 5 من أصل 28 كحد أقصى.} \\
\texttt{لقد جمعتما معاً 12 من أصل 16 عقوبة.} \\
\texttt{حرّك اللاعب الآخر كائناً على شبكته.} \\
\\ 
\\ 
\texttt{ما هو أمرك التالي؟} \\
            }
        }
    }
    & \\ \\

    \theutterance \stepcounter{utterance}  
    & \multicolumn{4}{p{0.6\linewidth}}{
        \cellcolor[rgb]{0.9,0.9,0.9}{
            \makecell[{{p{\linewidth}}}]{
                \texttt{\tiny{[P1$\rangle$GM]}}
                \texttt{حرّك: P, (2, 3)} \\
            }
        }
    }
    & & \\ \\

    \theutterance \stepcounter{utterance}  
    & & & \multicolumn{2}{p{0.3\linewidth}}{
        \cellcolor[rgb]{0.9,0.9,0.9}{
            \makecell[{{p{\linewidth}}}]{
                \texttt{\tiny{[GM$|$GM]}}
                \texttt{حرّك: P, (2, 3)} \\
            }
        }
    }
    & & \\ \\

    \theutterance \stepcounter{utterance}  
    & & & \multicolumn{2}{p{0.3\linewidth}}{
        \cellcolor[rgb]{0.9,0.9,0.9}{
            \makecell[{{p{\linewidth}}}]{
                \texttt{\tiny{[GM$|$GM]}}
                \texttt{P moved to (2, 3)} \\
            }
        }
    }
    & & \\ \\

    \theutterance \stepcounter{utterance}  
    & & \multicolumn{4}{p{0.6\linewidth}}{
        \cellcolor[rgb]{0.9,0.9,0.9}{
            \makecell[{{p{\linewidth}}}]{
                \texttt{\tiny{[GM$\rangle$P2]}}
                \texttt{تم تحريك 'S' إلى (5,7) بنجاح. تبدو شبكتك المحدّثة هكذا:} \\
\texttt{```} \\
\texttt{ 1234567} \\
\texttt{╔═╤══╤══╗} \\
\texttt{║◌│O◌│◌◌║ 1} \\
\texttt{║◌└┬─┴──╢ 2} \\
\texttt{║◌P│D◌◌◌║ 3} \\
\texttt{╟─┬┘◌◌┌─╢ 4} \\
\texttt{║◌│◌◌◌│A║ 5} \\
\texttt{╟─┼─┬─┤◌║ 6} \\
\texttt{║M│◌│S│T║ 7} \\
\texttt{╚═╧═╧═╧═╝} \\
\texttt{```} \\
\texttt{أنت تلعب حالياً الجولة 6 من أصل 28 كحد أقصى.} \\
\texttt{لقد جمعتما معاً 12 من أصل 16 عقوبة.} \\
\texttt{حرّك اللاعب الآخر كائناً على شبكته.} \\
\\ 
\\ 
\texttt{ما هو أمرك التالي؟} \\
            }
        }
    }
    & \\ \\

    \theutterance \stepcounter{utterance}  
    & & & \multicolumn{4}{p{0.6\linewidth}}{
        \cellcolor[rgb]{0.9,0.9,0.9}{
            \makecell[{{p{\linewidth}}}]{
                \texttt{\tiny{[GM$\langle$P2]}}
                \texttt{حرّك: S, (7, 6)} \\
            }
        }
    }
     \\ \\

    \theutterance \stepcounter{utterance}  
    & & & \multicolumn{2}{p{0.3\linewidth}}{
        \cellcolor[rgb]{0.9,0.9,0.9}{
            \makecell[{{p{\linewidth}}}]{
                \texttt{\tiny{[GM$|$GM]}}
                \texttt{حرّك: S, (7, 6)} \\
            }
        }
    }
    & & \\ \\

    \theutterance \stepcounter{utterance}  
    & & & \multicolumn{2}{p{0.3\linewidth}}{
        \cellcolor[rgb]{0.9,0.9,0.9}{
            \makecell[{{p{\linewidth}}}]{
                \texttt{\tiny{[GM$|$GM]}}
                \texttt{S moved to (7, 6)} \\
            }
        }
    }
    & & \\ \\

    \theutterance \stepcounter{utterance}  
    & & \multicolumn{4}{p{0.6\linewidth}}{
        \cellcolor[rgb]{0.9,0.9,0.9}{
            \makecell[{{p{\linewidth}}}]{
                \texttt{\tiny{[P1$\langle$GM]}}
                \texttt{تم تحريك 'P' إلى (2,3) بنجاح. تبدو شبكتك المحدّثة هكذا:} \\
\texttt{```} \\
\texttt{ 1234567} \\
\texttt{╔═╤══╤══╗} \\
\texttt{║◌│O◌│◌◌║ 1} \\
\texttt{║◌└┬─┴──╢ 2} \\
\texttt{║◌P│◌◌◌◌║ 3} \\
\texttt{╟─┬┘◌◌┌─╢ 4} \\
\texttt{║◌│AD◌│◌║ 5} \\
\texttt{╟─┼─┬─┤M║ 6} \\
\texttt{║◌│S│◌│T║ 7} \\
\texttt{╚═╧═╧═╧═╝} \\
\texttt{```} \\
\texttt{أنت تلعب حالياً الجولة 6 من أصل 28 كحد أقصى.} \\
\texttt{لقد جمعتما معاً 12 من أصل 16 عقوبة.} \\
\texttt{حرّك اللاعب الآخر كائناً على شبكته.} \\
\\ 
\\ 
\texttt{ما هو أمرك التالي؟} \\
            }
        }
    }
    & \\ \\

    \theutterance \stepcounter{utterance}  
    & \multicolumn{4}{p{0.6\linewidth}}{
        \cellcolor[rgb]{0.9,0.9,0.9}{
            \makecell[{{p{\linewidth}}}]{
                \texttt{\tiny{[P1$\rangle$GM]}}
                \texttt{حرّك: A, (7, 5)} \\
            }
        }
    }
    & & \\ \\

    \theutterance \stepcounter{utterance}  
    & & & \multicolumn{2}{p{0.3\linewidth}}{
        \cellcolor[rgb]{0.9,0.9,0.9}{
            \makecell[{{p{\linewidth}}}]{
                \texttt{\tiny{[GM$|$GM]}}
                \texttt{حرّك: A, (7, 5)} \\
            }
        }
    }
    & & \\ \\

    \theutterance \stepcounter{utterance}  
    & & & \multicolumn{2}{p{0.3\linewidth}}{
        \cellcolor[rgb]{0.9,0.9,0.9}{
            \makecell[{{p{\linewidth}}}]{
                \texttt{\tiny{[GM$|$GM]}}
                \texttt{A moved to (7, 5)} \\
            }
        }
    }
    & & \\ \\

    \theutterance \stepcounter{utterance}  
    & & \multicolumn{4}{p{0.6\linewidth}}{
        \cellcolor[rgb]{0.9,0.9,0.9}{
            \makecell[{{p{\linewidth}}}]{
                \texttt{\tiny{[GM$\rangle$P2]}}
                \texttt{تم تحريك 'S' إلى (7,6) بنجاح. تبدو شبكتك المحدّثة هكذا:} \\
\texttt{```} \\
\texttt{ 1234567} \\
\texttt{╔═╤══╤══╗} \\
\texttt{║◌│O◌│◌◌║ 1} \\
\texttt{║◌└┬─┴──╢ 2} \\
\texttt{║◌P│D◌◌◌║ 3} \\
\texttt{╟─┬┘◌◌┌─╢ 4} \\
\texttt{║◌│◌◌◌│A║ 5} \\
\texttt{╟─┼─┬─┤S║ 6} \\
\texttt{║M│◌│◌│T║ 7} \\
\texttt{╚═╧═╧═╧═╝} \\
\texttt{```} \\
\texttt{أنت تلعب حالياً الجولة 7 من أصل 28 كحد أقصى.} \\
\texttt{لقد جمعتما معاً 12 من أصل 16 عقوبة.} \\
\texttt{حرّك اللاعب الآخر كائناً على شبكته.} \\
\\ 
\\ 
\texttt{ما هو أمرك التالي؟} \\
            }
        }
    }
    & \\ \\

    \theutterance \stepcounter{utterance}  
    & & & \multicolumn{4}{p{0.6\linewidth}}{
        \cellcolor[rgb]{0.9,0.9,0.9}{
            \makecell[{{p{\linewidth}}}]{
                \texttt{\tiny{[GM$\langle$P2]}}
                \texttt{حرّك: S, (7, 6)} \\
            }
        }
    }
     \\ \\

    \theutterance \stepcounter{utterance}  
    & & & \multicolumn{2}{p{0.3\linewidth}}{
        \cellcolor[rgb]{0.9,0.9,0.9}{
            \makecell[{{p{\linewidth}}}]{
                \texttt{\tiny{[GM$|$GM]}}
                \texttt{حرّك: S, (7, 6)} \\
            }
        }
    }
    & & \\ \\

    \theutterance \stepcounter{utterance}  
    & & & \multicolumn{2}{p{0.3\linewidth}}{
        \cellcolor[rgb]{0.9,0.9,0.9}{
            \makecell[{{p{\linewidth}}}]{
                \texttt{\tiny{[GM$|$GM]}}
                \texttt{عقوبة: (7,6) ليست فارغة، بل تحتوي على 'S'.} \\
\\ 
\texttt{لقد جمعتما معاً 13 من أصل 16 عقوبة.} \\
\texttt{يرجى المحاولة مرة أخرى!} \\
            }
        }
    }
    & & \\ \\

    \theutterance \stepcounter{utterance}  
    & & \multicolumn{4}{p{0.6\linewidth}}{
        \cellcolor[rgb]{0.9,0.9,0.9}{
            \makecell[{{p{\linewidth}}}]{
                \texttt{\tiny{[GM$\rangle$P2]}}
                \texttt{عقوبة: (7,6) ليست فارغة، بل تحتوي على 'S'.} \\
\\ 
\texttt{لقد جمعتما معاً 13 من أصل 16 عقوبة.} \\
\texttt{يرجى المحاولة مرة أخرى!} \\
            }
        }
    }
    & \\ \\

    \theutterance \stepcounter{utterance}  
    & & & \multicolumn{4}{p{0.6\linewidth}}{
        \cellcolor[rgb]{0.9,0.9,0.9}{
            \makecell[{{p{\linewidth}}}]{
                \texttt{\tiny{[GM$\langle$P2]}}
                \texttt{حرّك: S, (6, 6)} \\
\\ 
\texttt{(ملاحظة للمدير: سأحاول تحريك S إلى (6,6) ولكنني أدركت أن (6,6) تحتوي على '┤'. سأحاول تحريك S إلى (5,6) بدلاً من ذلك إذا كانت فارغة، أو سأحاول تحريك كائن آخر إلى مكان فارغ تماماً مثل (2,1)).} \\
\\ 
\texttt{حرّك: S, (5, 6)} \\
            }
        }
    }
     \\ \\

    \theutterance \stepcounter{utterance}  
    & & & \multicolumn{2}{p{0.3\linewidth}}{
        \cellcolor[rgb]{0.9,0.9,0.9}{
            \makecell[{{p{\linewidth}}}]{
                \texttt{\tiny{[GM$|$GM]}}
                \texttt{حرّك: S, (6, 6)} \\
\\ 
\texttt{(ملاحظة للمدير: سأحاول تحريك S إلى (6,6) ولكنني أدركت أن (6,6) تحتوي على '┤'. سأحاول تحريك S إلى (5,6) بدلاً من ذلك إذا كانت فارغة، أو سأحاول تحريك كائن آخر إلى مكان فارغ تماماً مثل (2,1)).} \\
\\ 
\texttt{حرّك: S, (5, 6)} \\
            }
        }
    }
    & & \\ \\

    \theutterance \stepcounter{utterance}  
    & & & \multicolumn{2}{p{0.3\linewidth}}{
        \cellcolor[rgb]{0.9,0.9,0.9}{
            \makecell[{{p{\linewidth}}}]{
                \texttt{\tiny{[GM$|$GM]}}
                \texttt{Invalid format: head is not empty: 'حرّك: S, (6, 6)} \\
\\ 
\texttt{(ملاحظة للمدير: سأحاول تحريك S إلى (6,6) ولكنني أدركت أن (6,6) تحتوي على '┤'. سأحاول تحريك S إلى (5,6) بدلاً من ذلك إذا كانت فارغة، أو سأحاول تحريك كائن آخر إلى مكان فارغ تماماً مثل (2,1)).} \\
\\ 
\texttt{'} \\
            }
        }
    }
    & & \\ \\

    \theutterance \stepcounter{utterance}  
    & & \multicolumn{4}{p{0.6\linewidth}}{
        \cellcolor[rgb]{0.9,0.9,0.9}{
            \makecell[{{p{\linewidth}}}]{
                \texttt{\tiny{[GM$\rangle$P2]}}
                \texttt{عقوبة: يجب ألا تحتوي رسالتك على أي شيء قبل الأمر!} \\
\\ 
\texttt{تأكد من أن ردك يحتوي فقط على `قل: <رسالة>` أو `حرّك: <كائن>, (<X>, <Y>)`، ولا شيء غير ذلك!} \\
\\ 
\texttt{لقد جمعتما معاً 14 من أصل 16 عقوبة.} \\
\texttt{يرجى المحاولة مرة أخرى!} \\
            }
        }
    }
    & \\ \\

    \theutterance \stepcounter{utterance}  
    & & & \multicolumn{4}{p{0.6\linewidth}}{
        \cellcolor[rgb]{0.9,0.9,0.9}{
            \makecell[{{p{\linewidth}}}]{
                \texttt{\tiny{[GM$\langle$P2]}}
                \texttt{حرّك: S, (5, 6)} \\
            }
        }
    }
     \\ \\

    \theutterance \stepcounter{utterance}  
    & & & \multicolumn{2}{p{0.3\linewidth}}{
        \cellcolor[rgb]{0.9,0.9,0.9}{
            \makecell[{{p{\linewidth}}}]{
                \texttt{\tiny{[GM$|$GM]}}
                \texttt{حرّك: S, (5, 6)} \\
            }
        }
    }
    & & \\ \\

    \theutterance \stepcounter{utterance}  
    & & & \multicolumn{2}{p{0.3\linewidth}}{
        \cellcolor[rgb]{0.9,0.9,0.9}{
            \makecell[{{p{\linewidth}}}]{
                \texttt{\tiny{[GM$|$GM]}}
                \texttt{عقوبة: (5,6) ليست فارغة، بل تحتوي على '─'.} \\
\\ 
\texttt{لقد جمعتما معاً 15 من أصل 16 عقوبة.} \\
\texttt{يرجى المحاولة مرة أخرى!} \\
            }
        }
    }
    & & \\ \\

    \theutterance \stepcounter{utterance}  
    & & \multicolumn{4}{p{0.6\linewidth}}{
        \cellcolor[rgb]{0.9,0.9,0.9}{
            \makecell[{{p{\linewidth}}}]{
                \texttt{\tiny{[GM$\rangle$P2]}}
                \texttt{عقوبة: (5,6) ليست فارغة، بل تحتوي على '─'.} \\
\\ 
\texttt{لقد جمعتما معاً 15 من أصل 16 عقوبة.} \\
\texttt{يرجى المحاولة مرة أخرى!} \\
            }
        }
    }
    & \\ \\

    \theutterance \stepcounter{utterance}  
    & & & \multicolumn{4}{p{0.6\linewidth}}{
        \cellcolor[rgb]{0.9,0.9,0.9}{
            \makecell[{{p{\linewidth}}}]{
                \texttt{\tiny{[GM$\langle$P2]}}
                \texttt{حرّك: M, (2, 7)} \\
            }
        }
    }
     \\ \\

    \theutterance \stepcounter{utterance}  
    & & & \multicolumn{2}{p{0.3\linewidth}}{
        \cellcolor[rgb]{0.9,0.9,0.9}{
            \makecell[{{p{\linewidth}}}]{
                \texttt{\tiny{[GM$|$GM]}}
                \texttt{حرّك: M, (2, 7)} \\
            }
        }
    }
    & & \\ \\

    \theutterance \stepcounter{utterance}  
    & & & \multicolumn{2}{p{0.3\linewidth}}{
        \cellcolor[rgb]{0.9,0.9,0.9}{
            \makecell[{{p{\linewidth}}}]{
                \texttt{\tiny{[GM$|$GM]}}
                \texttt{عقوبة: (2,7) ليست فارغة، بل تحتوي على '│'.} \\
\\ 
\texttt{لقد جمعتما معاً 16 من أصل 16 عقوبة.} \\
\texttt{يرجى المحاولة مرة أخرى!} \\
            }
        }
    }
    & & \\ \\

    \theutterance \stepcounter{utterance}  
    & & \multicolumn{4}{p{0.6\linewidth}}{
        \cellcolor[rgb]{0.9,0.9,0.9}{
            \makecell[{{p{\linewidth}}}]{
                \texttt{\tiny{[GM$\rangle$P2]}}
                \texttt{عقوبة: (2,7) ليست فارغة، بل تحتوي على '│'.} \\
\\ 
\texttt{لقد جمعتما معاً 16 من أصل 16 عقوبة.} \\
\texttt{يرجى المحاولة مرة أخرى!} \\
            }
        }
    }
    & \\ \\

    \theutterance \stepcounter{utterance}  
    & & & \multicolumn{4}{p{0.6\linewidth}}{
        \cellcolor[rgb]{0.9,0.9,0.9}{
            \makecell[{{p{\linewidth}}}]{
                \texttt{\tiny{[GM$\langle$P2]}}
                \texttt{حرّك: M, (2, 1)} \\
            }
        }
    }
     \\ \\

    \theutterance \stepcounter{utterance}  
    & & & \multicolumn{2}{p{0.3\linewidth}}{
        \cellcolor[rgb]{0.9,0.9,0.9}{
            \makecell[{{p{\linewidth}}}]{
                \texttt{\tiny{[GM$|$GM]}}
                \texttt{حرّك: M, (2, 1)} \\
            }
        }
    }
    & & \\ \\

    \theutterance \stepcounter{utterance}  
    & & & \multicolumn{2}{p{0.3\linewidth}}{
        \cellcolor[rgb]{0.9,0.9,0.9}{
            \makecell[{{p{\linewidth}}}]{
                \texttt{\tiny{[GM$|$GM]}}
                \texttt{عقوبة: (2,1) ليست فارغة، بل تحتوي على '│'.} \\
\\ 
\texttt{لقد جمعتما معاً 17 من أصل 16 عقوبة.} \\
\texttt{يرجى المحاولة مرة أخرى!} \\
            }
        }
    }
    & & \\ \\

    \theutterance \stepcounter{utterance}  
    & & \multicolumn{4}{p{0.6\linewidth}}{
        \cellcolor[rgb]{0.9,0.9,0.9}{
            \makecell[{{p{\linewidth}}}]{
                \texttt{\tiny{[GM$\rangle$P2]}}
                \texttt{عقوبة: (2,1) ليست فارغة، بل تحتوي على '│'.} \\
\\ 
\texttt{لقد جمعتما معاً 17 من أصل 16 عقوبة.} \\
\texttt{يرجى المحاولة مرة أخرى!} \\
            }
        }
    }
    & \\ \\

    \theutterance \stepcounter{utterance}  
    & & & \multicolumn{4}{p{0.6\linewidth}}{
        \cellcolor[rgb]{0.9,0.9,0.9}{
            \makecell[{{p{\linewidth}}}]{
                \texttt{\tiny{[GM$\langle$P2]}}
                \texttt{حرّك: M, (2, 7)} \\
\\ 
\texttt{(ملاحظة للمدير: لقد تجاوزنا 16 عقوبة، ولكن سأحاول تنفيذ حركة صحيحة لإنقاذ الموقف إذا كان ذلك ممكناً. سأبحث عن خلية فارغة تماماً. بالنظر للشبكة، الخلية (2,1) تحتوي على '│'، والخلية (2,7) تحتوي على '│'. الخلية (2,2) تحتوي على '└'. الخلية (2,4) تحتوي على '┬'. الخلية (2,5) تحتوي على '│'. الخلية (2,6) تحتوي على '┼'. الخلية (3,2) تحتوي على '┬'. الخلية (3,4) تحتوي على '┘'. الخلية (3,5) تحتوي على '│'. الخلية (3,6) تحتوي على '─'. الخلية (3,7) تحتوي على '│'. الخلية (4,1) تحتوي على 'O'. الخلية (4,2) تحتوي على '◌'. الخلية (4,3) تحتوي على '◌'. الخلية (4,4) تحتوي على '◌'. الخلية (4,5) تحتوي على '◌'. الخلية (4,6) تحتوي على '─'. الخلية (4,7) تحتوي على '◌'. الخلية (5,1) تحتوي على '◌'. الخلية (5,2) تحتوي على '◌'. الخلية (5,3) تحتوي على '◌'. الخلية (5,4) تحتوي على '◌'. الخلية (5,5) تحتوي على '◌'. الخلية (5,6) تحتوي على '─'. الخلية (5,7) تحتوي على '◌'. الخلية (6,1) تحتوي على '◌'. الخلية (6,2) تحتوي على '◌'. الخلية (6,3) تحتوي على '◌'. الخلية (6,4) تحتوي على '◌'. الخلية (6,5) تحتوي على '◌'. الخلية (6,6) تحتوي على '┤'. الخلية (6,7) تحتوي على '│'. الخلية (7,1) تحتوي على '◌'. الخلية (7,2) تحتوي على '◌'. الخلية (7,3) تحتوي على '◌'. الخلية (7,4) تحتوي على '◌'. الخلية (7,5) تحتوي على 'A'. الخلية (7,6) تحتوي على 'S'. الخلية (7,7) تحتوي على 'T'.} \\
\\ 
\texttt{سأحاول تحريك M إلى (5, 1) فهي تبدو فارغة.)} \\
\\ 
\texttt{حرّك: M, (5, 1)} \\
            }
        }
    }
     \\ \\

    \theutterance \stepcounter{utterance}  
    & & & \multicolumn{2}{p{0.3\linewidth}}{
        \cellcolor[rgb]{0.9,0.9,0.9}{
            \makecell[{{p{\linewidth}}}]{
                \texttt{\tiny{[GM$|$GM]}}
                \texttt{حرّك: M, (2, 7)} \\
\\ 
\texttt{(ملاحظة للمدير: لقد تجاوزنا 16 عقوبة، ولكن سأحاول تنفيذ حركة صحيحة لإنقاذ الموقف إذا كان ذلك ممكناً. سأبحث عن خلية فارغة تماماً. بالنظر للشبكة، الخلية (2,1) تحتوي على '│'، والخلية (2,7) تحتوي على '│'. الخلية (2,2) تحتوي على '└'. الخلية (2,4) تحتوي على '┬'. الخلية (2,5) تحتوي على '│'. الخلية (2,6) تحتوي على '┼'. الخلية (3,2) تحتوي على '┬'. الخلية (3,4) تحتوي على '┘'. الخلية (3,5) تحتوي على '│'. الخلية (3,6) تحتوي على '─'. الخلية (3,7) تحتوي على '│'. الخلية (4,1) تحتوي على 'O'. الخلية (4,2) تحتوي على '◌'. الخلية (4,3) تحتوي على '◌'. الخلية (4,4) تحتوي على '◌'. الخلية (4,5) تحتوي على '◌'. الخلية (4,6) تحتوي على '─'. الخلية (4,7) تحتوي على '◌'. الخلية (5,1) تحتوي على '◌'. الخلية (5,2) تحتوي على '◌'. الخلية (5,3) تحتوي على '◌'. الخلية (5,4) تحتوي على '◌'. الخلية (5,5) تحتوي على '◌'. الخلية (5,6) تحتوي على '─'. الخلية (5,7) تحتوي على '◌'. الخلية (6,1) تحتوي على '◌'. الخلية (6,2) تحتوي على '◌'. الخلية (6,3) تحتوي على '◌'. الخلية (6,4) تحتوي على '◌'. الخلية (6,5) تحتوي على '◌'. الخلية (6,6) تحتوي على '┤'. الخلية (6,7) تحتوي على '│'. الخلية (7,1) تحتوي على '◌'. الخلية (7,2) تحتوي على '◌'. الخلية (7,3) تحتوي على '◌'. الخلية (7,4) تحتوي على '◌'. الخلية (7,5) تحتوي على 'A'. الخلية (7,6) تحتوي على 'S'. الخلية (7,7) تحتوي على 'T'.} \\
\\ 
\texttt{سأحاول تحريك M إلى (5, 1) فهي تبدو فارغة.)} \\
\\ 
\texttt{حرّك: M, (5, 1)} \\
            }
        }
    }
    & & \\ \\

    \theutterance \stepcounter{utterance}  
    & & & \multicolumn{2}{p{0.3\linewidth}}{
        \cellcolor[rgb]{0.9,0.9,0.9}{
            \makecell[{{p{\linewidth}}}]{
                \texttt{\tiny{[GM$|$GM]}}
                \texttt{Invalid format: head is not empty: 'حرّك: M, (2, 7)} \\
\\ 
\texttt{(ملاحظة للمدير: لقد تجاوزنا 16 عقوبة، ولكن سأحاول تنفيذ حركة صحيحة لإنقاذ الموقف إذا كان ذلك ممكناً. سأبحث عن خلية فارغة تماماً. بالنظر للشبكة، الخلية (2,1) تحتوي على '│'، والخلية (2,7) تحتوي على '│'. الخلية (2,2) تحتوي على '└'. الخلية (2,4) تحتوي على '┬'. الخلية (2,5) تحتوي على '│'. الخلية (2,6) تحتوي على '┼'. الخلية (3,2) تحتوي على '┬'. الخلية (3,4) تحتوي على '┘'. الخلية (3,5) تحتوي على '│'. الخلية (3,6) تحتوي على '─'. الخلية (3,7) تحتوي على '│'. الخلية (4,1) تحتوي على 'O'. الخلية (4,2) تحتوي على '◌'. الخلية (4,3) تحتوي على '◌'. الخلية (4,4) تحتوي على '◌'. الخلية (4,5) تحتوي على '◌'. الخلية (4,6) تحتوي على '─'. الخلية (4,7) تحتوي على '◌'. الخلية (5,1) تحتوي على '◌'. الخلية (5,2) تحتوي على '◌'. الخلية (5,3) تحتوي على '◌'. الخلية (5,4) تحتوي على '◌'. الخلية (5,5) تحتوي على '◌'. الخلية (5,6) تحتوي على '─'. الخلية (5,7) تحتوي على '◌'. الخلية (6,1) تحتوي على '◌'. الخلية (6,2) تحتوي على '◌'. الخلية (6,3) تحتوي على '◌'. الخلية (6,4) تحتوي على '◌'. الخلية (6,5) تحتوي على '◌'. الخلية (6,6) تحتوي على '┤'. الخلية (6,7) تحتوي على '│'. الخلية (7,1) تحتوي على '◌'. الخلية (7,2) تحتوي على '◌'. الخلية (7,3) تحتوي على '◌'. الخلية (7,4) تحتوي على '◌'. الخلية (7,5) تحتوي على 'A'. الخلية (7,6) تحتوي على 'S'. الخلية (7,7) تحتوي على 'T'.} \\
\\ 
\texttt{سأحاول تحريك M إلى (5, 1) فهي تبدو فارغة.)} \\
\\ 
\texttt{'} \\
            }
        }
    }
    & & \\ \\

    \theutterance \stepcounter{utterance}  
    & & & \multicolumn{2}{p{0.3\linewidth}}{
        \cellcolor[rgb]{0.9,0.9,0.9}{
            \makecell[{{p{\linewidth}}}]{
                \texttt{\tiny{[GM$|$GM]}}
                \texttt{Maximum number of penalties exceeded} \\
            }
        }
    }
    & & \\ \\

    \theutterance \stepcounter{utterance}  
    & & & \multicolumn{2}{p{0.3\linewidth}}{
        \cellcolor[rgb]{0.9,0.9,0.9}{
            \makecell[{{p{\linewidth}}}]{
                \texttt{\tiny{[GM$|$GM]}}
                \texttt{Initial states:} \\
\texttt{```} \\
\texttt{Player 1:} \\
\texttt{ 1234567} \\
\texttt{╔═╤══╤══╗} \\
\texttt{║◌│◌T│◌◌║ 1} \\
\texttt{║◌└┬─┴──╢ 2} \\
\texttt{║◌◌│O◌◌◌║ 3} \\
\texttt{╟─┬┘P◌┌─╢ 4} \\
\texttt{║◌│AD◌│◌║ 5} \\
\texttt{╟─┼─┬─┤M║ 6} \\
\texttt{║◌│S│◌│◌║ 7} \\
\texttt{╚═╧═╧═╧═╝} \\
\\ 
\\ 
\texttt{Player 2:} \\
\texttt{ 1234567} \\
\texttt{╔═╤══╤══╗} \\
\texttt{║◌│O◌│◌◌║ 1} \\
\texttt{║◌└┬─┴──╢ 2} \\
\texttt{║◌P│◌D◌◌║ 3} \\
\texttt{╟─┬┘◌◌┌─╢ 4} \\
\texttt{║◌│◌◌◌│A║ 5} \\
\texttt{╟─┼─┬─┤S║ 6} \\
\texttt{║M│◌│◌│T║ 7} \\
\texttt{╚═╧═╧═╧═╝} \\
\\ 
\texttt{```} \\
            }
        }
    }
    & & \\ \\

    \theutterance \stepcounter{utterance}  
    & & & \multicolumn{2}{p{0.3\linewidth}}{
        \cellcolor[rgb]{0.9,0.9,0.9}{
            \makecell[{{p{\linewidth}}}]{
                \texttt{\tiny{[GM$|$GM]}}
                \texttt{End states:} \\
\texttt{```} \\
\texttt{Player 1:} \\
\texttt{ 1234567} \\
\texttt{╔═╤══╤══╗} \\
\texttt{║◌│O◌│◌◌║ 1} \\
\texttt{║◌└┬─┴──╢ 2} \\
\texttt{║◌P│◌◌◌◌║ 3} \\
\texttt{╟─┬┘◌◌┌─╢ 4} \\
\texttt{║◌│◌D◌│A║ 5} \\
\texttt{╟─┼─┬─┤M║ 6} \\
\texttt{║◌│S│◌│T║ 7} \\
\texttt{╚═╧═╧═╧═╝} \\
\\ 
\\ 
\texttt{Player 2:} \\
\texttt{ 1234567} \\
\texttt{╔═╤══╤══╗} \\
\texttt{║◌│O◌│◌◌║ 1} \\
\texttt{║◌└┬─┴──╢ 2} \\
\texttt{║◌P│D◌◌◌║ 3} \\
\texttt{╟─┬┘◌◌┌─╢ 4} \\
\texttt{║◌│◌◌◌│A║ 5} \\
\texttt{╟─┼─┬─┤S║ 6} \\
\texttt{║M│◌│◌│T║ 7} \\
\texttt{╚═╧═╧═╧═╝} \\
\\ 
\texttt{```} \\
            }
        }
    }
    & & \\ \\

    \theutterance \stepcounter{utterance}  
    & & & \multicolumn{2}{p{0.3\linewidth}}{
        \cellcolor[rgb]{0.9,0.9,0.9}{
            \makecell[{{p{\linewidth}}}]{
                \texttt{\tiny{[GM$|$GM]}}
                \texttt{* success: False} \\
\texttt{* lose: False} \\
\texttt{* aborted: True} \\
\texttt{{-}{-}{-}{-}{-}{-}{-}} \\
\texttt{* Shifts: 6.00} \\
\texttt{* Max Shifts: 12.00} \\
\texttt{* Min Shifts: 6.00} \\
\texttt{* End Distance Sum: 12.21} \\
\texttt{* Init Distance Sum: 27.62} \\
\texttt{* Expected Distance Sum: 29.33} \\
\texttt{* Penalties: 18.00} \\
\texttt{* Max Penalties: 16.00} \\
\texttt{* Rounds: 6.00} \\
\texttt{* Max Rounds: 28.00} \\
\texttt{* Object Count: 7.00} \\
            }
        }
    }
    & & \\ \\

\end{supertabular}
}

\end{document}
